\documentclass[12pt]{article}
\usepackage[utf8]{inputenc}
\usepackage[T1]{fontenc}
\usepackage{hyperref}
\title{Makhuwa-Saka Language Resources for K-12 Curriculum Development}
\author{Compiled by Claude Sonnet 5.5 (claude-sonnet-5-5) for Andr\'as Kornai}
\date{}
\begin{document}
\maketitle

\section{Introduction}
Makhuwa-Saka (ISO 639-3 xsq; Glottocode makh1266) is a Bantu variety of the Makhuwa cluster in Mozambique. Alternate names include Esaka, Saka, Esaaka, Isaanga, Ishanga and Sanga. Wikipedia gives about 210,000 speakers of eSaka (Ethnologue, 2006), with limited mutual intelligibility with other Makhuwa varieties. Joshua Project gives 477,000 to 491,000 people. Makhuwa peoples live in Zambezia, Nampula and Cabo Delgado provinces. The searches found only general profiles and no school materials specific to Makhuwa-Saka.

\section{Primary Educational Resources}
No primers, textbooks or school programmes specific to Makhuwa-Saka were found. Materials in other Makhuwa varieties may exist, but they were not found or verified here.

\section{Linguistic and Reference Materials}
\begin{itemize}
\item Wikipedia, ``Makhuwa language'', which covers the cluster and lists the varieties. \url{https://en.wikipedia.org/wiki/Makhuwa_language}
\item Wikitongues, Makhuwa-Saka (xsq) language page. \url{https://wikitongues.org/languages/xsq/}
\end{itemize}

\section{Religious and Community Resources}
\begin{itemize}
\item Joshua Project, Makhuwa-Saka in Mozambique. It marks the language as written and lists Bible portions. \url{https://joshuaproject.net/people_groups/19512/MZ}
\item Joshua Project, language page for xsq. \url{https://joshuaproject.net/languages/xsq}
\end{itemize}

\section{Digital and Audio Resources}
Global Recordings Network has audio in the related Esaka variety, but none was reported for Makhuwa-Saka itself. \url{https://glorecnet-website.k.globalrecordings.net/en/language/25571}

\section{Curriculum Development Context}
No evidence was found of Makhuwa-Saka in schools. Any curriculum would build on the written Bible portions and on work in other Makhuwa varieties. It would need an orthography and teaching materials developed with local speakers, and a decision on how Saka relates to standard Makhuwa given the limited mutual intelligibility. The speaker counts found are rough.
\end{document}